\addcontentsline{toc}{subsubsection}{Solution 58. The Four-Step Prime Labyrinth}

\begin{dlsummary}[title={\bfseries Solution 58. The Four-Step Prime Labyrinth},breakable]

\begin{proof}
Write \(f^{\circ j}\) for \(j\) successive applications of \(f\).

\paragraph{1. Force equality by taking reciprocals.}

Positivity and multiplicativity give
\[
f(1)=1,\qquad f(1/x)=\frac1{f(x)}.
\]
The same reciprocal identity holds for every iterate. Applying the given inequality at \(1/x\) therefore gives the reverse bound, so
\[
\boxed{f^{\circ4}(x)=x^4.} \tag{1}
\]

For every prime \(p\), we must have \(f(p)>1\); otherwise all subsequent iterates would equal \(1\), contradicting (1).

Also, multiplicativity and integer preservation imply
\[
a\mid b\quad\Longrightarrow\quad f(a)\mid f(b)
\]
for positive integers \(a,b\).

\paragraph{2. Prove that each prime maps to a prime power.}

Fix a prime \(p\). Choose primes successively so that
\[
q\mid f(p),\qquad r\mid f(q),\qquad s\mid f(r).
\]
These choices are possible because the image of every prime exceeds \(1\).

By repeated divisibility,
\[
f(s)\mid f^{\circ4}(p)=p^4.
\]
Since \(f(s)>1\), we have \(p\mid f(s)\). Thus we have a closed divisibility chain:
\[
p\longrightarrow q\longrightarrow r\longrightarrow s\longrightarrow p,
\]
where an arrow \(u\longrightarrow v\) means \(v\mid f(u)\).

Now let \(q'\) be \textbf{any} prime divisor of \(f(p)\). Restarting at \(q\), we obtain
\[
q\longrightarrow r\longrightarrow s\longrightarrow p\longrightarrow q'.
\]
Hence
\[
q'\mid f^{\circ4}(q)=q^4,
\]
so \(q'=q\).

Therefore \(f(p)\) has only one prime divisor:
\[
\boxed{f(p)=\sigma(p)^{a_p},}
\]
where \(\sigma(p)\) is prime and \(a_p\) is a positive integer.

\paragraph{3. Classify the prime cycles.}

Applying \(f\) four times,
\[
f^{\circ4}(p)
=
\sigma^4(p)^{
a_p a_{\sigma(p)}a_{\sigma^2(p)}a_{\sigma^3(p)}
}.
\]
Comparing with \(p^4\),
\[
\boxed{
\sigma^4(p)=p,\qquad
a_p a_{\sigma(p)}a_{\sigma^2(p)}a_{\sigma^3(p)}=4.
} \tag{2}
\]

Thus \(\sigma\) is a permutation whose cycles have lengths \(1,2,\) or \(4\).

\begin{itemize}
\item **Length \(1\) is impossible:** It would require \(a_p^4=4\).
\item **Length \(2\):** If \(p\leftrightarrow q\), then \((a_pa_q)^2=4\), so \(a_pa_q=2\).
\item **Length \(4\):** The four exponents must have product \(4\).
\end{itemize}

\paragraph{Complete classification.}

Partition the primes into disjoint directed cycles of the following types:

| Cycle | Assignments |
|---|---|
| Two primes \(p,q\) | \(f(p)=q,\quad f(q)=p^2\), after choosing an orientation |
| Four primes \(p_1,p_2,p_3,p_4\) | \(f(p_i)=p_{i+1}^{a_i}\), with indices taken cyclically and \(a_1a_2a_3a_4=4\) |

For a four-cycle, the exponent pattern consists either of \textbf{one \(4\) and three \(1\)s}, or \textbf{two \(2\)s and two \(1\)s}, in any cyclic arrangement.

Extend these prime assignments multiplicatively:
\[
f\!\left(\prod_p p^{e_p}\right)
=\prod_p f(p)^{e_p},
\qquad e_p\in\mathbb Z.
\]
These are exactly all solutions.

\textbf{Verification.} The construction preserves positive integers and is multiplicative. On a two-cycle, four iterations produce exponent \((1\cdot2)^2=4\); on a four-cycle, they produce exponent \(a_1a_2a_3a_4=4\). Therefore \(f^{\circ4}(x)=x^4\) for every positive rational \(x\).
\end{proof}

\paragraph{Takeaways.}
\emph{No takeaways were provided in the source draft.}

\end{dlsummary}
